\documentclass[a4paper]{article}
% Español (México) con XeLaTeX: fontspec para fuentes Unicode, polyglossia
% para la división silábica y los nombres del documento en la variante mexicana.
\usepackage{fontspec}
\usepackage{polyglossia}
\setdefaultlanguage[variant=mexican]{spanish}
% Los nombres chinos van como «pinyin (original)»: xeCJK da a los caracteres
% Han su propia fuente; sin ella XeLaTeX los omite con un aviso.
% FandolSong no cubre algunos caracteres de nombres (U+73FA); WenQuanYi Zen Hei sí.
\usepackage[AutoFallBack=true]{xeCJK}
\setCJKmainfont{FandolSong-Regular.otf}
\setCJKfallbackfamilyfont{\CJKrmdefault}{WenQuanYi Zen Hei}
% polyglossia no traduce los nombres de listings.
\AtBeginDocument{\providecommand{\lstlistingname}{}\renewcommand{\lstlistingname}{Listado}%
  \providecommand{\lstlistlistingname}{}\renewcommand{\lstlistlistingname}{Índice de listados}}
\usepackage{amsmath, amssymb}
\usepackage{graphicx}
\usepackage[margin=2.35cm]{geometry}
\usepackage[most]{tcolorbox}
\usepackage{booktabs}
\usepackage{tabularx}
\usepackage{longtable}
\usepackage{array}
\usepackage{float}
\usepackage{xcolor}
\usepackage{hyperref}
\usepackage{fancyhdr}
\hypersetup{colorlinks=true, linkcolor=blue!55!black, urlcolor=blue!60!black}
\graphicspath{{./}{figures/}}
\setlength{\parskip}{0.55em}
\setlength{\parindent}{2em}
\setlength{\emergencystretch}{3em}
\renewcommand{\arraystretch}{1.18}
\newtcolorbox{knowledgebox}[1]{enhanced, colback=blue!5!white, colframe=blue!70!black, colbacktitle=blue!70!black, coltitle=white, fonttitle=\bfseries, title={#1}, sharp corners, breakable}
\newtcolorbox{importantbox}[1]{enhanced, colback=yellow!10!white, colframe=yellow!75!black, colbacktitle=yellow!75!black, coltitle=black, fonttitle=\bfseries, title={#1}, sharp corners, breakable}
\newtcolorbox{warningbox}[1]{enhanced, colback=red!5!white, colframe=red!70!black, colbacktitle=red!70!black, coltitle=white, fonttitle=\bfseries, title={#1}, sharp corners, breakable}
\pagestyle{fancy}
\fancyhf{}
\fancyfoot[C]{\thepage}
\fancyfoot[R]{\footnotesize\color{gray}\href{https://github.com/hqhq1025/ai-course-notes}{AI Course Notes}}
\renewcommand{\headrulewidth}{0pt}
\renewcommand{\footrulewidth}{0pt}
\newcommand{\videourl}{https://www.youtube.com/watch?v=owjTOT14bG0}
\newcommand{\repourl}{https://github.com/hqhq1025/ai-course-notes}
\begin{document}
\begin{titlepage}
\centering
{\Large Notas de panorama técnico\par}
\vspace{1.0cm}
{\huge\bfseries Panorama de los datos: el petróleo de la nueva era, la pirámide, la fijación de precios y la Recipe\par}
\vspace{0.6cm}
{\Large Zhang Xiaojun conversa con Xie Chen (谢晨)\par}
\vspace{0.4cm}
{\large Elaborado a partir del video público de Zhang Xiaojun Podcast, una transcripción hecha con GPU y diagramas conceptuales propios\par}
\vspace{0.3cm}
{\large 2026-03-30\par}
\vspace{0.9cm}
\includegraphics[width=0.88\textwidth,height=0.42\textheight,keepaspectratio]{cover.jpg}\par
\vfill
\begin{tcolorbox}[width=0.92\textwidth, colback=black!2!white, colframe=black!55, sharp corners]
\textbf{Título del video}: 134. «Panorama de los datos». Conversación con Xie Chen: el petróleo de la nueva era, historia, mapa del sector, pirámide de datos, fijación de precios y Recipe\par
\textbf{Canal del video}: Zhang Xiaojun Podcast\par
\textbf{Duración del video}: 02:38:23\par
\textbf{Enlace del video}: \href{\videourl}{\nolinkurl{\videourl}}\par
\textbf{Nota sobre los materiales}: YouTube no tiene subtítulos disponibles; el texto se elaboró a partir de una transcripción con faster-whisper large-v3 en CUDA, los capítulos del video, la portada y diagramas didácticos propios. Las tomas fijas de la entrevista no se reutilizan en el texto.\par
\textbf{Página del proyecto}: \href{\repourl}{\nolinkurl{\repourl}}\par
\end{tcolorbox}
\end{titlepage}
\tableofcontents
\newpage
\section{Introducción: por qué los datos vuelven a ser el eje de la industria}

Suele decirse que la inteligencia artificial se impulsa con tres motores: los datos, la capacidad de cómputo y los algoritmos. En los últimos años, la capacidad de cómputo y la arquitectura de los modelos han sido más visibles; pero ahora que los modelos de lenguaje de gran tamaño se acercan al «muro de datos» y la robótica enfrenta un «desierto de datos», los datos vuelven a ser la variable que determina el panorama de la industria. El valor de esta entrevista con Xie Chen (谢晨) es que no habla de un solo conjunto de datos, sino que descompone la industria de los datos en fuentes, ciclo cerrado, pirámide, fijación de precios, Recipe y mapa del sector.

\begin{importantbox}{Tesis central de este episodio}
Los datos no son una materia prima estática, sino un ciclo cerrado continuo. Quien obtenga a menor costo fallas reales, escenarios simulados, retroalimentación verificable y recetas reutilizables podrá entrenar más rápido sistemas inteligentes que se puedan llevar a la práctica.
\end{importantbox}

\begin{knowledgebox}{Nota sobre la estrategia visual}
Este video es una entrevista con encuadre fijo, sin slides didácticas, pizarrón ni demostraciones de producto. Conforme al estándar de este repositorio para pódcasts, el cuerpo del texto no repite fotogramas de los participantes; la portada sirve para identificar la fuente, y el contenido didáctico se presenta mediante la pirámide de datos, el ciclo cerrado y el mapa de la industria.
\end{knowledgebox}

\subsection{Resumen de la sección}

Este episodio es un panorama de la industria de los datos: el problema de los LLM es el rendimiento marginal decreciente de los datos de alta calidad; el de la robótica, que los datos del mundo físico son escasos y costosos. Ambos necesitan un ciclo cerrado de datos y un sistema de evaluación.
\section{La pirámide de datos: datos reales, de simulación y de internet}

La mayor diferencia entre los datos de robótica y los de los LLM es el anclaje físico. Un LLM puede obtener de internet cantidades enormes de texto, código y material multimodal; un robot necesita acciones, estados, entorno y retroalimentación, y recolectarlos en el mundo real es costoso y lento. Por eso Xie Chen (谢晨) enfatiza que los datos de robótica se parecen más a una pirámide: en la cima están los datos reales recolectados en el propio dispositivo, en el nivel intermedio los datos de simulación o sintéticos, y en la base los datos de internet y de personas.

\begin{figure}[H]
\centering
\includegraphics[width=0.94\textwidth]{data-pyramid.png}
\caption{La pirámide de datos de robótica: datos reales, datos de simulación y datos de internet o de personas.}
\end{figure}

\begin{knowledgebox}{Lectura de la figura: qué resuelve cada uno de los tres niveles de datos}
Los datos reales de la cima son los más cercanos al despliegue, pero son escasos y costosos; los datos de simulación del nivel intermedio son controlables y escalables, adecuados para generar escenarios de cola larga; los datos de internet o de personas de la base tienen la mayor escala y el menor costo, pero su anclaje físico es débil. Una estrategia de datos para robótica tiene que combinar los tres niveles.
\end{knowledgebox}

\subsection{Por qué los datos de robótica son un desierto}

Los robots no son como la conducción autónoma, donde millones de vehículos envían todos los días datos reales de las calles. Los robots de propósito general todavía no se han desplegado a gran escala, la teleoperación es costosa, los escenarios de las tareas están dispersos y los cuerpos de los robots difieren mucho entre sí. Sin escala de despliegue en los dispositivos no surge de forma natural un motor de datos como el de Tesla. Por eso, para la robótica, la simulación y los datos sintéticos son una condición previa, no solo un complemento.

\begin{warningbox}{La simulación no es un juguete, pero tampoco una panacea}
La simulación permite generar escenarios y evaluaciones a escala, pero existe la sim-to-real gap. Los datos reales son escasos, la simulación es imprecisa y los datos de internet tienen un anclaje débil: los tres deben combinarse, no sustituirse entre sí.
\end{warningbox}

\subsection{Resumen de la sección}

La contradicción central de los datos de robótica es esta: los datos más reales son los más escasos, y los más abundantes son los menos reales. La pirámide de datos es el marco básico para resolver esta contradicción.

\section{Motor de datos: el ciclo cerrado de datos que va del fracaso al éxito}

Xie Chen (谢晨) dice que los datos más eficaces son los de «primero fallar y luego tener éxito». Esta afirmación es importante. Las trayectorias exitosas le indican al modelo qué debe hacer; las trayectorias que fallan y luego se corrigen le indican por qué se equivocó, cómo recuperarse y cómo aprender del error. Tanto la conducción autónoma como los robots necesitan este ciclo cerrado: despliegue, recolección, minería, entrenamiento, evaluación y nuevo despliegue.

\begin{figure}[H]
\centering
\includegraphics[width=0.94\textwidth]{data-engine-loop.png}
\caption{Ciclo cerrado de datos / Data Engine: despliegue, recolección, minería, entrenamiento, evaluación, simulación.}
\end{figure}

\begin{knowledgebox}{Lectura de la figura: lo esencial del Data Engine es la continuidad}
En la figura, cada ronda de despliegue trae de vuelta retroalimentación real; los casos de fracaso se extraen y se convierten en datos; después del entrenamiento se validan mediante evaluación en simulación y, por último, se vuelve a desplegar. La conclusión que sustenta es esta: los datos no son una compra única, sino el resultado de que el sistema opere de forma continua.
\end{knowledgebox}

\subsection{Lo que enseña el motor de datos de Tesla}

El motor de datos de Tesla se alimenta de una gran cantidad de vehículos en el lado del dispositivo. Los autos operan en el mundo real, generan escenarios de cola larga y retroalimentación de conducción, que regresan a la nube para el entrenamiento y actualizan las capacidades en el lado del dispositivo. Esta lógica no se cumple del todo en los robots, porque la escala de robots en el lado del dispositivo es insuficiente. Por ello, la industria robótica debe buscar un nuevo motor de datos: una combinación de simulación, síntesis, teleoperación, minería de casos de fracaso, benchmark y despliegue real a pequeña escala.

\begin{importantbox}{La esencia del motor de datos}
Un motor de datos no consiste en «tener muchos datos», sino en contar con un ciclo capaz de detectar errores de manera continua, generar datos complementarios, validar las mejoras y volver a desplegar.
\end{importantbox}

\subsection{Resumen de la sección}

Los datos de alta calidad provienen de un ciclo cerrado. Tanto los robots como los LLM necesitan retroalimentación, pero la retroalimentación de los robots es más costosa y depende más del entorno y de la evaluación.
\section{Mapa de tipos de datos: LLM, modelos del mundo, VLA y robots}

En la entrevista se mencionan varias veces los modelos de lenguaje de gran tamaño, los modelos del mundo, VLA, el cerebro de los robots y la IA del mundo físico. Estos conceptos se confunden con facilidad. Pueden entenderse ubicándolos en un mapa de tipos de datos: los datos de texto y código son adecuados para los LLM, los datos multimodales aportan percepción, las trayectorias de robots aportan acciones, los datos de simulación aportan escenarios controlables y los pares fracaso-éxito aportan retroalimentación con alta densidad de información.

\begin{figure}[H]
\centering
\includegraphics[width=0.96\textwidth]{data-taxonomy.png}
\caption{Mapa de tipos de datos: texto, código, multimodal, trayectorias de robots, simulación y pares fracaso-éxito.}
\end{figure}

\begin{knowledgebox}{Lectura de la figura: cada modelo necesita datos distintos}
Los LLM dependen del texto y del código; los modelos del mundo necesitan datos relacionados con la predicción física; VLA necesita visión, lenguaje y acción; las políticas de robots necesitan trayectorias y retroalimentación obtenidas en el propio dispositivo. La escasez de datos no significa que no existan datos en absoluto, sino que faltan datos que correspondan a la capacidad que se busca.
\end{knowledgebox}

\subsection{Términos clave: modelos y datos}

\noindent\begin{tabularx}{\textwidth}{p{0.22\textwidth}p{0.32\textwidth}X}
\toprule
Término & Problema que resuelve & Datos que requiere \\
\midrule
LLM & Comprensión del lenguaje, razonamiento, código y tareas del mundo digital & Texto, código, corpus multimodales, datos de retroalimentación. \\
World Model & Predecir el estado del mundo físico y sus cambios futuros & Video, simulación, interacción física, datos 3D y espaciales. \\
VLA & Vision-Language-Action: conecta percepción, lenguaje y acción & Visión, instrucciones en lenguaje, trayectorias de acción, retroalimentación del entorno. \\
Physical AI & Sistemas de IA capaces de actuar en el mundo físico & Datos del cuerpo del robot, simulación, despliegue real, evaluación. \\
Behavior Benchmark & Conjunto de evaluación de largo horizonte orientado a tareas corporizadas & Entornos de simulación, definición de tareas, criterios de éxito. \\
\bottomrule
\end{tabularx}

\subsection{Relación simbiótica}

Los equipos de modelos grandes, de modelos del mundo y de VLA no están aislados. VLA puede usar un LLM base, un modelo del mundo puede funcionar como cerebro en la nube y VLA puede devolver al modelo del mundo la retroalimentación de sus acciones. Si los sistemas de evaluación convergen poco a poco, la frontera entre los modelos del mundo y VLA también podría estrecharse.

\begin{warningbox}{No confundas los modelos del mundo con VLA}
Los modelos del mundo se enfocan más en comprender y predecir el mundo físico; VLA se enfoca más en actuar en el mundo físico. Ambos coexistirán en simbiosis, pero tienen objetivos distintos, datos distintos y también formas de evaluación distintas.
\end{warningbox}

\subsection{Resumen de la sección}

El mapa de datos determina el mapa de modelos. Los LLM, los modelos del mundo, VLA y las políticas de robots necesitan tipos de datos distintos, y también se van conectando mediante la evaluación y la retroalimentación.
\section{Fijación de precios de los datos: los datos no se venden por GB}

A los datos se les llama el petróleo de la nueva era, pero esta metáfora puede inducir a error. El petróleo se vende por cantidad física, mientras que el valor de los datos depende de su autenticidad, escasez, verificabilidad, reutilización y beneficio marginal. Un mismo 1GB de datos puede ser una colección de páginas web repetidas o un conjunto de casos de falla clave de la cola larga; su valor es completamente distinto.

\begin{figure}[H]
\centering
\includegraphics[width=0.92\textwidth]{data-pricing.png}
\caption{Las cinco dimensiones de la fijación de precios de los datos: autenticidad, escasez, verificabilidad, reutilización y beneficio marginal.}
\end{figure}

\begin{knowledgebox}{Lectura de la figura: el precio de los datos proviene de su contribución marginal}
Los cinco factores de la figura determinan en conjunto el valor de los datos. Cuanto más auténticos, escasos y verificables sean, y cuanto mejor se transfieran, más valen; sin embargo, cuantos más datos del mismo tipo haya, menor es su beneficio marginal. Lo esencial de la industria de datos no es acumular volumen, sino encontrar datos de alto beneficio marginal.
\end{knowledgebox}

\subsection{Por qué los datos de fallas son caros}

Los datos de fallas son caros porque le indican al modelo dónde están sus límites. Muchas trayectorias exitosas pueden parecerse entre sí, mientras que las trayectorias fallidas suelen exponer defectos del sistema, escenarios de la cola larga y estrategias de recuperación. Para la conducción autónoma, los robots y los Agent, los pares de falla y corrección tienen más valor de entrenamiento que las simples demostraciones exitosas.

\begin{importantbox}{Los datos más efectivos: primero la falla, después el éxito}
Los datos que primero fallan y después tienen éxito contienen el error, el diagnóstico, la corrección y el resultado. Se acercan más al proceso de aprendizaje que los datos de puro éxito y también son más adecuados para entrenar la capacidad de recuperación.
\end{importantbox}

\subsection{Resumen de la sección}

El precio de los datos debe fijarse según su contribución marginal a la mejora de capacidades, no según su volumen. Los datos de alto valor suelen ser auténticos, escasos y verificables, y son capaces de exponer los límites del modelo.

\section{Data Recipe: la receta de datos es la verdadera barrera}

Con la misma fuente de datos, distintas formas de limpieza, filtrado, proporción de mezcla, orden de entrenamiento y retroalimentación producen capacidades distintas en el modelo. La recipe es la barrera menos visible de la industria de datos. No consiste en «comprar un montón de datos», sino en decidir qué datos entran al entrenamiento, en qué proporción, en qué momento y cómo se evalúan.

\begin{figure}[H]
\centering
\includegraphics[width=0.92\textwidth]{data-recipe.png}
\caption{Data Recipe: fuentes, filtrado, proporción de mezcla, currículo y retroalimentación de la evaluación.}
\end{figure}

\begin{knowledgebox}{Lectura de la figura: por qué la recipe es más difícil de copiar que la materia prima}
En el centro de la figura está la Data Recipe y a su alrededor aparecen las fuentes, el filtrado, la proporción de mezcla, el currículo y la evaluación. La materia prima se puede comprar; la receta surge de la experimentación, los fracasos, la evaluación y el conocimiento del dominio. La recipe determina si los datos se convierten realmente en capacidad.
\end{knowledgebox}

\subsection{Términos clave: conceptos relacionados con la recipe}

\noindent\begin{tabularx}{\textwidth}{p{0.22\textwidth}p{0.32\textwidth}X}
\toprule
Término & Problema que resuelve & Significado en este episodio \\
\midrule
Filtering & Eliminar datos de baja calidad, duplicados, dañinos o irrelevantes & Determina la pureza de la señal de entrenamiento. \\
Mixture & Proporción entre datos de distintas fuentes & Determina la distribución de capacidades y las preferencias. \\
Curriculum & Orden en que los datos entran al entrenamiento & Determina qué aprende primero el modelo y qué aprende después. \\
Eval Feedback & Usar los resultados de la evaluación para deducir qué cambiar en los datos & Hace que la recipe sea iterable. \\
Synthetic Data & Datos generados por un modelo o por simulación & Compensa la escasez de datos reales, pero exige controlar el sesgo de distribución. \\
\bottomrule
\end{tabularx}

\subsection{Resumen de la sección}

La Data Recipe es el convertidor que transforma los datos en capacidad. Lo que de verdad resulta difícil de copiar en la industria de datos no suele ser la fuente de datos en sí, sino la receta y el ciclo cerrado de retroalimentación.
\section{Panorama de la industria: quién compite por el ciclo cerrado de datos de la robótica y la IA física}

En la entrevista se mencionan varios grupos de actores: los equipos de modelos grandes, los equipos de modelos del mundo, los equipos de VLA, los fabricantes de robots y las empresas de simulación y datos sintéticos. No forman una simple cadena de proveedores y clientes, sino una relación de simbiosis. Quien controle el ciclo cerrado de datos tiene más probabilidades de controlar el cerebro del robot.

\begin{figure}[H]
\centering
\includegraphics[width=0.95\textwidth]{data-industry-landscape.png}
\caption{Panorama de la industria de datos: modelos grandes, modelos del mundo, VLA, fabricantes de robots y empresas de datos de simulación.}
\end{figure}

\begin{knowledgebox}{Lectura de la figura: por qué es una red simbiótica}
Los equipos de modelos grandes aportan el modelo base y la infraestructura; los equipos de modelos del mundo se ocupan de la predicción física; los equipos de VLA conectan con las acciones; los fabricantes de robots tienen los robots reales; las empresas de simulación aportan escenarios a escala. Ninguna de las partes puede completar por sí sola todo el ciclo cerrado de datos.
\end{knowledgebox}

\subsection{Fabricantes de robots frente a empresas del cerebro robótico}

En la conducción autónoma, las empresas que tienen una flota de vehículos cuentan de forma natural con datos del lado del dispositivo; en la robótica todavía no existe un despliegue a gran escala en dispositivos, por lo que los fabricantes de robots no necesariamente tienen una ventaja de datos por defecto. Si una empresa del cerebro robótico controla la simulación y la infraestructura, también puede desarrollar capacidades primero; si un fabricante de robots logra formar un ciclo cerrado con datos reales, también será muy relevante.

\begin{warningbox}{La experiencia de la conducción autónoma no se traslada directamente a la robótica}
La conducción autónoma cuenta con entornos viales relativamente estandarizados y con flotas de gran escala; en la robótica, los cuerpos de los robots, las tareas y los escenarios están mucho más dispersos. El motor de datos debe rediseñarse; no se puede copiar la ruta de FSD.
\end{warningbox}

\subsection{Resumen de la sección}

La industria de datos para robótica será una red simbiótica. Los modelos, los modelos del mundo, los VLA, los fabricantes de robots y las empresas de datos de simulación pueden convertirse en nodos clave; lo decisivo es quién logra formar un ciclo cerrado.
\section{El punto final: ¿dejarán de importar los datos?}

Al final de la entrevista se pregunta: ¿llegará el día en que el problema de los datos deje de importar por completo? Xie Chen (谢晨) responde que cada vez está más convencido de que, cuanto más fuerte es la inteligencia, más hambre tiene de conocimiento y de datos. Cuanto más capaz es una persona, más quiere aprender, y con la IA podría ocurrir lo mismo. El desenlace quizá no sea que «los datos desaparezcan», sino pasar de aprender de datos externos a aprender de forma continua en simulación y en entornos propios.

\subsection{La evaluación es el cuello de botella}

En el caso de los robots, Xie Chen considera que el problema más importante podría ser escalar la evaluación. Sin una evaluación a escala no se sabe si el modelo de verdad se volvió más inteligente. Los modelos de lenguaje grandes y los Agent también necesitan métricas de evaluación de nivel superior, porque cuanto más fuerte es el modelo, más fácil es que los benchmark anteriores dejen de servir.

\begin{importantbox}{La evaluación también es un problema de datos}
El conjunto de evaluación no es un juez externo, sino una fuente de retroalimentación para el entrenamiento y la investigación. Sin una buena evaluación, el ciclo cerrado de datos no puede saber hacia dónde avanzar.
\end{importantbox}

\subsection{Autoaprendizaje y el desenlace en simulación}

Si la IA puede, dentro de un entorno de simulación, fijarse objetivos, intentar, fallar, corregir y volver a intentar, la fábrica de datos no desaparecerá, sino que cambiará de forma: pasará de recolectar datos humanos a construir entornos, objetivos y retroalimentación para que la IA aprenda en ellos de forma continua.

\begin{knowledgebox}{De la data factory a la environment factory}
Cuando el rendimiento marginal de los datos externos disminuye, lo que de verdad importa podría ser construir entornos de los que se pueda aprender. La simulación, la evaluación, la reward y la retroalimentación de los fallos se convertirán en la nueva fábrica de datos.
\end{knowledgebox}

\subsection{Resumen de la sección}

Los datos no desaparecerán: pasarán de ser un corpus estático a ser entornos dinámicos. Cuanto más fuerte es la inteligencia, más necesita retroalimentación de mayor calidad y entornos de aprendizaje más complejos.
\section{Diferencias entre el muro de datos de los LLM y el desierto de datos de la robótica}

Aunque en ambos casos se dice «faltan datos», lo que les falta a los LLM y a los robots no es lo mismo. El problema de los LLM se parece más a un rendimiento marginal decreciente del corpus público de alta calidad: las páginas web, los libros, el código, los artículos científicos y los datos de preguntas y respuestas ya se han usado a gran escala, y seguir ampliándolos tropieza con cuellos de botella de repetición, contaminación, derechos de autor, calidad y evaluación. A la robótica, en cambio, le faltan desde el principio datos de interacción real: si no hay suficientes robots ejecutando tareas en entornos reales, no existe un flujo estable de datos desde los dispositivos.

\subsection{Dos tipos de escasez de datos}

\noindent\begin{longtable}{p{0.22\textwidth}p{0.35\textwidth}p{0.34\textwidth}}
\toprule
Objeto & Tipo de escasez de datos & Vías de solución habituales \\
\midrule
LLM & Rendimiento marginal decreciente del corpus de alta calidad; los benchmark se vuelven más difíciles & Mejor filtrado, mejores datos sintéticos, retroalimentación de orden superior, entornos de agent. \\
Robótica & Pocas trayectorias de acción reales, pocos escenarios de cola larga, recolección costosa & Simulación, teleoperación, despliegue real a pequeña escala, minería de datos de fallos. \\
Modelos del mundo & Datos insuficientes de predicción física y de dinámica espacial & Video, 3D, simulación, tareas de predicción multimodal. \\
VLA & Instrucciones, visión, acciones y retroalimentación difíciles de unificar & Trayectorias capturadas en el dispositivo, tareas de simulación, anotación en lenguaje, action head. \\
\bottomrule
\end{longtable}

\begin{importantbox}{El «muro de datos» no es un concepto unificado}
El muro de datos de los LLM es un problema de rendimiento marginal y de calidad; el desierto de datos de la robótica es un problema de interacción física y de escala de despliegue. Confundir estos dos conceptos conduce a estrategias de datos equivocadas.
\end{importantbox}

\subsection{Resumen de la sección}

Tanto a los LLM como a los robots les faltan datos, pero de manera distinta. Los primeros necesitan datos de mayor calidad y retroalimentación de orden superior; los segundos necesitan interacción física escalable y un ciclo cerrado de simulación.

\section{Métodos de recolección de datos para robots: real, teleoperación, simulación y síntesis}

La recolección de datos para robots puede seguir varias vías, y cada una tiene su costo y su sesgo. El despliegue real es el más cercano al producto, pero no alcanza la escala necesaria; la teleoperación permite obtener demostraciones de acciones, pero su costo en mano de obra es alto; la simulación permite escalar los escenarios de cola larga, pero tiene un sim-to-real gap; los videos de internet ofrecen gran escala, pero carecen de acciones y de retroalimentación. Un sistema de datos de alto nivel tiene que combinar estas fuentes.

\subsection{Comparación de los métodos de recolección}

\noindent\begin{longtable}{p{0.22\textwidth}p{0.30\textwidth}p{0.39\textwidth}}
\toprule
Método & Ventajas & Costos \\
\midrule
Despliegue real & Es el más cercano a las tareas y al hardware reales & Escala reducida, riesgo alto, pocos casos de cola larga. \\
Teleoperación & Demostraciones de acciones claras y alto control & Mano de obra costosa; la distribución de operadores es limitada. \\
Simulación & Puede generarse a escala y cubrir escenarios peligrosos o poco frecuentes & Brecha con la realidad y costo del modelado físico. \\
Datos sintéticos & Permiten completar con rapidez tareas y anotaciones de lenguaje & Sesgo de distribución y riesgo de alucinaciones. \\
Videos de internet & Escala enorme y cobertura de escenarios muy diversos & Carecen de acciones, de estados y de reward. \\
\bottomrule
\end{longtable}

\begin{knowledgebox}{Por qué los datos de fallos son los más valiosos}
Las trayectorias fallidas contienen condiciones límite e información de recuperación. Para un robot, las demostraciones exitosas le dicen al modelo «cómo hacerlo», mientras que las correcciones tras un fallo le dicen «dónde se equivoca» y «cómo recuperarse». Este tipo de datos tiene más valor marginal que repetir demostraciones exitosas.
\end{knowledgebox}

\subsection{Resumen de la sección}

La recolección de datos para robots no sigue una sola vía, sino que combina varias fuentes. El despliegue real, la teleoperación, la simulación, la síntesis y los videos de internet resuelven, cada uno, problemas distintos.
\section{Calidad de los datos de simulación: cómo abordar la sim-to-real gap}

La simulación es la capa intermedia más importante de la pirámide de datos para robots. Permite generar escenarios a escala, controlar variables, reproducir fallas y construir evaluaciones. Sin embargo, los datos de simulación solo tienen valor cuando la diferencia con el mundo real está bajo control. De lo contrario, el modelo aprende los atajos de la simulación y no las capacidades que necesita en la realidad.

\subsection{Las cuatro dimensiones de la calidad de la simulación}

\noindent\begin{tabularx}{\textwidth}{p{0.22\textwidth}p{0.34\textwidth}X}
\toprule
Dimensión & Requisito & Manifestación de la falla \\
\midrule
Coherencia física & Contacto, fricción, gravedad y dinámica creíbles & Las acciones fallan en la realidad. \\
Realismo visual & Materiales, iluminación, oclusiones y ruido de cámara razonables & El modelo de percepción se sobreajusta a las imágenes de la simulación. \\
Diversidad de tareas & Cubrir la cola larga, las anomalías y las combinaciones complejas & Solo resuelve escenarios estándar. \\
Transferibilidad de la evaluación & La evaluación en simulación predice el desempeño real & Las tablas de clasificación en simulación no corresponden con el despliegue real. \\
\bottomrule
\end{tabularx}

\begin{warningbox}{Más realismo no siempre es mejor simulación}
Perseguir en exceso el realismo visual puede tener un costo alto y un beneficio bajo. La simulación debe servir a los objetivos de entrenamiento y evaluación: hay que distinguir con claridad qué variables influyen en la política y cuáles solo afectan la apariencia.
\end{warningbox}

\subsection{Resumen de la sección}

La calidad de los datos de simulación no depende de «si se ve parecido o no», sino de si mejora el desempeño en tareas reales y ayuda a evaluar las capacidades reales.
\section{La evaluación es el punto de llegada y también el de partida}

Al final de la entrevista se enfatiza que tanto la robótica como los grandes modelos enfrentan un cuello de botella en la evaluación. Sin una buena evaluación no es posible saber si el modelo de verdad se volvió más inteligente; sin una evaluación que pueda escalarse no es posible cerrar el ciclo de datos. La evaluación no es un informe que se redacta al terminar el entrenamiento, sino el punto de partida que determina la siguiente ronda de recolección de datos y la dirección del entrenamiento.

\subsection{Los tres papeles de la evaluación}

\noindent\begin{longtable}{p{0.22\textwidth}p{0.34\textwidth}p{0.35\textwidth}}
\toprule
Papel & Función & Ejemplos \\
\midrule
Árbitro & Juzgar si el modelo actual es más capaz & benchmark, tasa de éxito de extremo a extremo, tasa de aprobación en tareas reales. \\
Minero de datos & Exponer los fallos y los escenarios de cola larga & Los casos de fallo regresan de forma automática al conjunto de entrenamiento. \\
Volante de dirección & Decidir la siguiente dirección de entrenamiento y de recolección & Se agregan datos del tipo de tarea en el que la carencia es mayor. \\
\bottomrule
\end{longtable}

\begin{importantbox}{La evaluación también es un activo de datos}
Una evaluación de alta calidad define «qué cuenta como progreso». A medida que los modelos se vuelven más capaces, la evaluación misma se convierte en un dato escaso, porque su diseño requiere personas que entiendan mejor la tarea, los fallos y las necesidades futuras.
\end{importantbox}

\subsection{Resumen de la sección}

El ciclo de datos debe tener la evaluación como centro. Sin evaluación no hay dirección; sin dirección, los datos solo se acumulan y no se convierten en capacidad.
\section{Cadena de suministro de datos: el proceso industrial que va de la fuente de datos a la capacidad}

Si se concibe los datos como una industria, hace falta una cadena de suministro: recolección, licenciamiento, limpieza, anotación, síntesis, mezcla, entrenamiento, evaluación y retroalimentación. Cada eslabón influye en la capacidad final del modelo. Cuanto más larga es la cadena de suministro de datos, más necesarias se vuelven la estandarización, la auditoría y el control de calidad.

\subsection{Tabla de la cadena de suministro de datos}

\noindent\begin{longtable}{p{0.20\textwidth}p{0.34\textwidth}p{0.37\textwidth}}
\toprule
Eslabón & Tarea & Riesgo de calidad \\
\midrule
Recolección & Obtener datos de páginas web, robots, simulación y comportamiento de usuarios & Sesgo de origen, derechos de autor, privacidad. \\
Limpieza & Eliminar duplicados y ruido, filtrar muestras de baja calidad & Borrar por error casos de cola larga de alto valor. \\
Anotación & Dar estructura a tareas, acciones, estados y retroalimentación & Anotaciones inconsistentes, costo alto. \\
Síntesis & Usar modelos o simulación para cubrir escenarios escasos & Desplazamiento de la distribución, colapso de modos. \\
Mezcla & Definir la proporción de los datos y el orden de entrenamiento & Capacidades desequilibradas, olvido. \\
Evaluación & Revisar las capacidades del modelo y los tipos de falla & Métricas distorsionadas, sobreajuste al benchmark. \\
Retorno & Devolver las fallas y los escenarios nuevos al conjunto de datos & Ciclo cerrado lento, retroalimentación imprecisa. \\
\bottomrule
\end{longtable}

\begin{knowledgebox}{La barrera de entrada de la cadena de suministro de datos}
La barrera de entrada de los datos no consiste solo en «tengo datos», sino también en «sé cómo convertir los datos en capacidad». La limpieza, la proporción de la mezcla, la evaluación y el retorno suelen ser más difíciles de copiar que los datos en bruto.
\end{knowledgebox}

\subsection{Resumen de la sección}

La industrialización del sector de datos implica que entre la fuente de datos y la capacidad existe una cadena de suministro completa. Si cualquier eslabón es débil, la capacidad final del modelo disminuye.
\section{Conexión con los episodios anteriores: Agent, AI for Math y datos}

EP139 trata la historia de la tecnología de los Agent; EP138, el posentrenamiento de los Agent y la reasignación de la capacidad de cómputo; EP137, la verificación formal en AI for Math; y EP134 aporta la perspectiva de los datos, que es la base de todo lo anterior. Vistos en conjunto, los cuatro episodios tienen un punto en común: la siguiente etapa de la IA no dependerá solo de tener más parámetros, sino de mejores entornos, retroalimentación, evaluaciones y ciclos cerrados de datos.

\begin{knowledgebox}{Hilo conductor de los cuatro episodios}
Los Agent de EP139 necesitan retroalimentación del entorno; el posentrenamiento de EP138 necesita rollout y reward; Lean, en EP137, ofrece una verificación estricta; y el motor de datos de EP134 explica de dónde sale esa retroalimentación. En esencia, todos abordan la misma pregunta: «cómo lograr que la inteligencia aprenda de forma continua».
\end{knowledgebox}

\subsection{Resumen de la sección}

El panorama de los datos no es un tema aislado. Explica por qué los Agent, el posentrenamiento, AI for Math y la robótica necesitan entornos de retroalimentación de alta calidad.
\section{Gobernanza de datos: derechos de autor, privacidad y control de calidad}

En cuanto la industria de datos alcanza escala, los problemas de gobernanza se vuelven ineludibles. Los datos para LLM implican derechos de autor, privacidad y transparencia sobre su procedencia; los datos para robots implican la recolección en entornos reales, los escenarios de los usuarios, los registros de sensores y la responsabilidad en materia de seguridad; los datos de simulación, por su parte, implican supuestos físicos y sesgos de evaluación. Sin gobernanza, cuantos más datos haya, mayor será el riesgo.

\subsection{El triángulo de la gobernanza}

\noindent\begin{longtable}{p{0.22\textwidth}p{0.34\textwidth}p{0.35\textwidth}}
\toprule
Dimensión & Pregunta que debe responderse & Consecuencia de fallar \\
\midrule
Derechos de autor y licencias & ¿Los datos tienen una procedencia legal y un alcance de uso definido? & Controversias sobre los datos de entrenamiento y riesgos comerciales posteriores. \\
Privacidad y seguridad & ¿Los datos contienen información personal, ubicaciones, hogares, fábricas o secretos comerciales? & Filtración de información sensible e imposibilidad de entrar en entornos empresariales. \\
Calidad y trazabilidad & ¿De dónde provienen los datos, cómo se limpian y cómo influyen en el modelo? & No es posible localizar las fallas del modelo y la Recipe no es reproducible. \\
\bottomrule
\end{longtable}

\begin{warningbox}{La gobernanza de datos no es un trámite legal de último momento}
Si la procedencia, la limpieza y las licencias de los datos no son trazables, cuanto más capaz sea el modelo, mayor será el riesgo. Esto ocurre sobre todo en los robots y en los entornos empresariales, porque los datos suelen contener espacios reales, equipos reales y comportamientos reales de los usuarios.
\end{warningbox}

\subsection{Las cuatro preguntas del control de calidad}

Antes de que un conjunto de datos entre en el entrenamiento, hay que hacerse al menos cuatro preguntas: si representa la tarea objetivo; si contiene las capacidades que le faltan al modelo; si tiene etiquetas o retroalimentación confiables; si introducirá atajos erróneos. En el caso de los datos de simulación, además hay que preguntar: ¿la política aprendida en la simulación puede transferirse al mundo real?

\begin{importantbox}{El control de calidad es requisito previo de la Recipe}
Sin control de calidad, la Recipe no es más que una proporción de mezcla elegida casi por superstición; con control de calidad, la Recipe puede convertirse en un método de ingeniería reproducible.
\end{importantbox}

\subsection{Resumen de la sección}

La gobernanza de datos determina si los datos pueden usarse a largo plazo. Los derechos de autor, la privacidad, la trazabilidad y el control de calidad son el umbral que la industria de datos debe cruzar para pasar de la «recolección» a la «infraestructura».
\section{Lista práctica de Recipe: cómo convertir los datos en capacidades}

Data Recipe no es un eslogan, sino un proceso experimental. Si un equipo afirma tener una ventaja en datos, al menos debe poder responder: de dónde vienen los datos, cómo se filtran, cómo se mezclan, cómo entran al entrenamiento, cómo se evalúan, cómo regresan los fallos y cómo se ajusta la siguiente ronda.

\subsection{Lista práctica de siete pasos}

\noindent\begin{longtable}{p{0.16\textwidth}p{0.36\textwidth}p{0.39\textwidth}}
\toprule
Paso & Acción & Punto de verificación \\
\midrule
1. Definir la capacidad & Precisar qué capacidad de tarea se quiere mejorar & ¿Es percepción, planificación, acción, recuperación o generalización? \\
2. Buscar fuentes de datos & Reales, simulados, sintéticos, de internet, retroalimentación humana & ¿Los datos cubren la capacidad objetivo? \\
3. Limpiar y filtrar & Eliminar duplicados, ruido, contenido irrelevante y contaminación & ¿Se eliminaron por error casos de cola larga y casos de fallo? \\
4. Fijar las proporciones & Determinar la proporción de cada fuente y de cada nivel de dificultad & ¿Provoca un desequilibrio entre capacidades? \\
5. Currículo de entrenamiento & Decidir qué se aprende primero y qué después & ¿Respeta la progresión de las capacidades? \\
6. Evaluación y retroalimentación & Verificar con benchmark, simulación o tareas reales & ¿Las métricas se relacionan con el despliegue real? \\
7. Ciclo cerrado de retorno & Devolver los fallos, la cola larga y las tareas nuevas al conjunto de datos & ¿Los datos de la siguiente ronda tienen mayor rendimiento marginal? \\
\bottomrule
\end{longtable}

\begin{knowledgebox}{Lo que se puede replicar de una Recipe y lo que no}
El proceso se puede replicar, pero la receta concreta es difícil de replicar, porque surge de la interacción entre el modelo, la tarea, el hardware, la evaluación y la experiencia del equipo; si cambia la capacidad objetivo, quizá haya que rehacer la receta.
\end{knowledgebox}

\subsection{Resumen de la sección}

Para convertir los datos en capacidades hay que partir de la capacidad objetivo, no de los datos que se tienen a la mano. Recipe es ingeniería de datos orientada a capacidades.
\section{Compra y producción propia de datos: ¿comprar datos o crearlos?}

En la industria, las opciones habituales son comprar datos, recolectar datos y crear datos. Comprar datos es rápido, pero lleva a la homogeneidad; construir datos propios crea barreras profundas, pero es costoso; los datos sintéticos o de simulación son escalables, pero conllevan un riesgo de falta de realismo. La elección cambia según la etapa: en la validación temprana se puede comprar o sintetizar, pero al llegar a las capacidades centrales es indispensable construir un ciclo cerrado propio.

\subsection{Comparación de las tres vías}

\noindent\begin{tabularx}{\textwidth}{p{0.20\textwidth}p{0.34\textwidth}X}
\toprule
Vía & Etapa adecuada & Riesgos \\
\midrule
Comprar datos & Arranque rápido, complemento de capacidades generales & Homogeneidad, licencias poco claras, rendimiento marginal bajo. \\
Recolectar datos & Desarrollo de capacidades en tareas reales & Costo alto, plazos largos, privacidad y seguridad complejas. \\
Crear datos & Escenarios de cola larga, peligrosos o escasos & Desplazamiento de la distribución, evaluaciones distorsionadas, sobreajuste a la simulación. \\
\bottomrule
\end{tabularx}

\begin{importantbox}{Las barreras de largo plazo provienen de un ciclo cerrado propio}
Los datos que se pueden comprar difícilmente se convierten en una ventaja defensiva de largo plazo. La verdadera barrera proviene de la combinación de tareas propias, retroalimentación propia, evaluación propia y Recipe propia.
\end{importantbox}

\subsection{Resumen de la sección}

La estrategia de datos no consiste en elegir entre comprar o crear. A corto plazo se puede comprar y sintetizar; a largo plazo es indispensable formar un ciclo cerrado de datos y un sistema de evaluación propios.
\section{Términos clave: índice de palabras clave de este episodio}

\begin{longtable}{p{0.24\textwidth}p{0.32\textwidth}p{0.35\textwidth}}
\toprule
Término & Explicación en una frase & Función en este episodio \\
\midrule
Pirámide de datos & Estructura jerárquica de datos reales, simulados y de internet o humanos & Explica la estrategia de datos para robots. \\
Data Engine & Ciclo cerrado de despliegue, recolección, minería, entrenamiento y evaluación & Mecanismo central de la industria de datos. \\
Synthetic Data & Datos sintéticos o generados por simulación & Compensan la escasez de datos reales de robots. \\
World Model & Modelo que predice el estado y el futuro del mundo físico & Coexiste en simbiosis con los VLA. \\
VLA & Modelo Vision-Language-Action & Forma importante del cerebro de un robot. \\
Sim-to-real Gap & Brecha entre la simulación y el mundo real & Problema que los datos de simulación deben resolver. \\
Data Recipe & Estrategia de limpieza, proporción, orden y retroalimentación de los datos & Ahí reside la barrera competitiva de los datos. \\
Behavior Benchmark & Conjunto de evaluación de tareas corporizadas & Punto de convergencia entre la evaluación de robots y la de modelos del mundo. \\
\bottomrule
\end{longtable}

\subsection{Resumen de la sección}

Todas las palabras clave de este episodio apuntan a una misma conclusión: la industria de datos no consiste en vender datos, sino en construir un ciclo cerrado de datos sostenible, un sistema de evaluación y recetas de entrenamiento.

\section{Síntesis y ampliación}

\subsection{Conclusiones centrales}

\begin{enumerate}
\item Los LLM enfrentan rendimientos marginales decrecientes de los datos; la robótica enfrenta un desierto de datos reales.
\item Los datos para robótica deben entenderse como una pirámide: datos reales recopilados en el dispositivo, datos de simulación o sintéticos y datos de internet o de personas.
\item El núcleo del Data Engine es un ciclo continuo: detectar fallas, completar los datos, entrenar, evaluar y volver a desplegar.
\item El valor de los datos no se mide en GB, sino por su autenticidad, su escasez, su verificabilidad y su rendimiento marginal.
\item La Data Recipe es una barrera de entrada implícita de la industria de datos.
\item El panorama de la industria robótica es una red simbiótica de empresas de modelos grandes, de modelos del mundo, de VLA, de cuerpos robóticos y de datos de simulación.
\item El destino final de los datos no es que desaparezcan, sino el paso de los datos estáticos a los entornos dinámicos y al autoaprendizaje.
\end{enumerate}

\subsection{Preguntas abiertas}

\begin{itemize}
\item Dentro de la pirámide de datos para robótica, ¿pueden los datos de simulación superar de verdad la sim-to-real gap?
\item ¿Quién controlará el ciclo cerrado de datos de la robótica: las empresas de cuerpos robóticos, las de modelos grandes o las de datos de simulación?
\item ¿Puede la Data Recipe convertirse en una barrera central, como lo es la arquitectura del modelo?
\item Cuando se fortalezca el autoaprendizaje de la IA, ¿la data factory se convertirá en una environment factory?
\end{itemize}

\subsection{Lecturas adicionales}

\begin{itemize}
\item Tesla FSD data engine: para entender el ciclo cerrado de datos recopilados en el dispositivo.
\item Artículos sobre Behavior benchmark, VLA y world model: para entender la relación entre la evaluación de robots y los modelos del mundo.
\item Materiales sobre datos sintéticos, simulación, domain randomization y sim-to-real: para entender los retos centrales de los datos para robótica.
\item EP137 AI for Math: para entender cómo un entorno verificable se convierte en una fuente de datos de alta calidad.
\end{itemize}

\begin{importantbox}{Juicio final}
El futuro de los datos no son «más rastreadores web», sino «mejores entornos». Quien logre construir entornos verificables, escalables y capaces de generar retroalimentación de fallas tomará la iniciativa en la siguiente etapa del entrenamiento de la inteligencia.
\end{importantbox}

\end{document}
